%
% Praxissemester Bericht
% Tobias Birkle
% Vorlage: Sven Hodapp



% -----------
% 1. Präambel
% -----------


% Allgemeine Einstellungen
% ------------------------
\documentclass[
	pdftex,%              PDFTex verwenden da wir ausschliesslich ein PDF erzeugen.
	a4paper,%             Wir verwenden A4 Papier.
	oneside,%             Einseitiger Druck.
	12pt,%                Grosse Schrift, besser geeignet für A4.
	halfparskip,%         Halbe Zeile Abstand zwischen Absätzen.
	%chapterprefix,%       Kapitel mit 'Kapitel' anschreiben.
	headsepline,%         Linie nach Kopfzeile.
	footsepline,%         Linie vor Fusszeile.
	bibtotocnumbered,%    Literaturverzeichnis im Inhaltsverzeichnis nummeriert einfügen.
	idxtotoc%             Index ins Inhaltsverzeichnis einfügen.
]{article}

\usepackage[a4paper, includeheadfoot, bindingoffset=0cm, 
left=2cm, right=2cm, top=2cm, bottom=2 cm]{geometry}

\usepackage{fancyhdr}
\pagestyle{fancy} 
\renewcommand{\headrulewidth}{0pt}


\usepackage[utf8]{inputenc}
%\usepackage[german]{babel}   % deutsche Silbentrennung
%\selectlanguage{german}   % damit Table Of Contents Inhaltsverzeichnis genannt wird

\usepackage{geometry}   % Seitenränder einstellbar
\usepackage{textcomp}   % Sonderzeichen, wie Eurosymbol


% Absätze mit Leerzeile ohne Einrückung
\newenvironment{absaetze}
{\addtolength{\parskip}{\baselineskip}\parindent 0pt}
{}


% Bilder, Farben, farbige Tabellen
% --------------------------------
\usepackage{graphicx, color, colortbl}
\usepackage{array}       % Erweiterte Tabelleneigenschaften.
%\usepackage{floatflt}   % Bild kann von Text umflossen werden.



% Palatino Schrift
% ----------------
%\usepackage[T1]{fontenc}
%\usepackage[osf]{mathpazo}   % osf aktiviert Mediävalziffern/Minuskelziffern



% Sonstige Pakete
% ---------------
%\usepackage{anysize}   % Seitenränder verändern
%\usepackage{setspace}   % 1.5em Zeilenabstand \begin{onehalfspacing}
%\usepackage{bibgerm}   % Anzeigestil des Literaturverzeichnis (gerabbrv)
\usepackage{listings}  \lstset{numbers=left, numberstyle=\tiny, numbersep=20pt} % Programmcode einfügen
\usepackage{subfigure}


\renewcommand{\baselinestretch}{1.1}\normalsize
% PDF Eigenschaften
% -----------------
\usepackage
[
	colorlinks=false,
	bookmarks = true,
	pdftitle={Room surveillance app},
	pdfsubject={Documentation},
	pdfkeywords={Documentation, Ubicom, Room surveillance},
	pdfauthor={Benjamin Kugler, Tobias Birkle},
	urlcolor=blue,
	pdfstartview=FitH
]{hyperref}



\usepackage{float}

% --------------------
% 2. Dokumenten Anfang
% --------------------

\begin{document}


% Deckblatt
% ---------
\begin{absaetze}
\begin{titlepage}

	\includegraphics[width=1.00\textwidth]{images/HTWG_logo.PNG}

	\vspace*{1cm}
	\begin{center}
		\Huge
		\vspace{2cm}
		\textbf{Area surveillance system} \\
		\vspace{1cm}
		\large
		
		\vspace{1cm}
		\textbf{Benjamin Kugler, 286346} \\
		\textbf{Tobias Birkle, 286081} \\
		\today\\
		\vspace{2cm}
		
		
		Prof. Dr. rer. nat. Ralf E.D. Seepold\\
		

		
	\end{center}
	\normalsize
	\vfill
\end{titlepage}

\section*{Abstract}
The area surveillance system observes areas in your building. Whenever the system is active, it will take a picture if a movement is detected. The movements will be detected with a passive infra-red-sensor (PIR). The PIR is connected to an arduino board. The arduino and the camera are connected to a LAN. If the PIR is triggered, the arduino gets an image from the camera over the common gateway interface(CGI). The image will be uploaded to a webserver. The mobile phone app receives a notification message via the google cloud messaging server(GCM). Afterwards the user can download the image to his mobile phone.



\end{absaetze}

\begin{absaetze}
\renewcommand{\thepage}{\arabic{page}}
\fancyfoot[CE, CO]{}

\fancyfoot[RE, RO]{\thepage}
\fancyhead[]{}


\renewcommand{\thepage}{\Roman{page}}
%\addcontentsline{toc}{section}{Kurzfassung}
\end{absaetze}
% Inhaltsverzeichnis anzeigen
% ---------------------------
\setcounter{page}{1}
\tableofcontents
%\newpage
%\listoffigures % Abbildungsverzeichnis
%\listoftables




% ---------
% 3. Inhalt
% ---------

\begin{absaetze}
\section[State of the art - Team]{State of the art}
As of now, there are Cameras available to buy, which can do everything our area surveillance app can do. But they cost much more and the source code isn’t available. So you can’t add new features easily. For example the camera we used, the Lupus LE950 can detect faces and you can configure it so it will only observe a selected area in the picture. There are also system which can handle more cameras at once to observe more than one area.

\section[Goals and requirements]{Goals and requirements}
We wanted to create something similar to the Lupus LE950. A system which can take images if you walk through a sensor and which can upload the image to the internet/mobile phone. A simple surveillance system. The Lupus LE950 is very expensive and we wanted to create a system which costs less money than the Lupus and where we can tinker a little bit. The requirement was, that we have to use the arduino uno for the task. The arduino is limited in its processing power so we couldn’t implement features such face detection. We wanted to test the arduino, how fast can it process the image and how big can the images be. 

\section[Use Case - Team]{Use case}

\begin{figure}[H]
	\centering
	\includegraphics[width=0.9\textwidth]{images/usecase.png}
	\caption[Use-case diagram]{This image shows the use-case diagram of the area surveillance system}
	\label{usecase}
\end{figure}


As we described beforehand in the abstract topic you can see in the use-case diagram \ref{usecase} how the behaviour of our area surveillance system works.
Android App:
\begin{itemize}
\item receive notification: The app receives a notification from the gcm over our webserver. The notification contains the image URL.
\item download images: With the provided URL the user can download the image.
\item store data: After the image was downloaded, it will be stored on the mobile device
\item show gallery: If the user wants, he can review all the older pictures he downloaded earlier.
\end{itemize}

Webserver:
\begin{itemize}
\item notifiy android: The webserver receives a notification from the arduino to notify the android app.
\item store data: The webserver stores the images, received from the arduino, on its data store.
\end{itemize}

Surveillance System:
\begin{itemize}
\item retrieve sensor data: after the PIR sensor gets triggered, he sends a signal to the arduino to take a picture
\item take picture: if the sensor was triggered, the camera takes a picture.
\item receive image: the arduino will receive an image over the CGI script.
\item store data: after the image is taken, the arduino saves it on the sd card.
\item upload data: the saved image will be uploaded to the ftp server.
\item notify webserver: the arduino sends a notification to the webserver and triggers the php script.

\end{itemize}


Other use cases as described in the abstract could be a door camera system. The system gets activated if someone rings the bell or is near the door. So older people or wheelchair drivers can see whos there. We are aware that most older people doesn’t have mobile phones so we could show the image on a tablet or a big screen in the house.

Another use case could be for wildlife watching. You could mount the camera and the IR sensor on a tree or another surface and it will take pictures of the animals who cross the sensor.



\section[System description - Team]{System description}
In the next paragraph we will explain the system architecture and the components we used in the project.

\subsection{Overview}
\begin{figure}[H]
	\centering
	\includegraphics[width=0.9\textwidth]{images/overview.png}
	\caption[System overview]{This image shows the an overview of the system}
	\label{overview}
\end{figure}

In the image \ref{overview} you can see the system and its components. 
There are two parts. One part in the local area network connected with ethernet and the other part in the internet. We are using an Arduino UNO and a passive infrared sensor to detect movements. To get the images we use the LUPUS LE950 webcam. Also we need a web server to store the images in the internet and a mobile phone to download them.

\subsection{Data flow}
\begin{figure}[H]
	\centering
	\includegraphics[width=0.9\textwidth]{images/dataflow.png}
	\caption[System data flow]{This image shows the data flow of the system}
	\label{dataflow}
\end{figure}

In the figure \ref{dataflow} you can see the data flow of the system.
Firstly the Arduino UNO senses movement with the PIR-sensor. If the the sensor tells the arduino that it detected movement, the Arduino sends a request to the camera, gets an image and stores it on the Arduino’s SD-card.
%In the instruction manual there are few different CGI commands, so we asked the company if %they could provide us with more information about the camera and more CGI commands. We %received a list with commands but they were notvery usefull. So it  was more clear, that we %needed some more open source kind of stuff.

Then the Arduino sends the image to the web server via an FTP connection. After the image was uploaded to the web server, a php script renames the picture with a timestamp and sends a notification to the google cloud messaging server. Everyone who is registered in the gcm server will receive a notification on their mobile phones. So our app doesn’t need to poll all the time if there is new data available. This saves energy and therefore the battery can last longer. Within the notification the url for the download is also transmitted to the mobile phone. The image is on a ftp server and is easy to download.

\subsection{UML}

\section[Development environment - Team]{Development environment}
For the develompement of the area surveillance app, we used different tools and technologies to enhance the developement process.

\subsection{Tools}
\begin{itemize}
\item Arduino IDE. Easy to use IDE to upload the software to the Arduino UNO.
\item Eclipse with adt plugin. Eclipse is an easy to use IDE for Android developement
\item GIT. Version control system.
\item Wireshark. We needed wireshark to evaluate the communication from the camera over the cgi script.
\end{itemize}

\subsection{Technologies}

\textbf{CGI - Common Gateway Interface} \\
The common gateway interface is a method to generate dynamic content on web pages or web applications. They are normally written in a scripting language like shell. In this project the camera provides us a set of cgi scripts to get images or adjust the camera.

\textbf{PHP - PHP: Hypertext Preprocessor} \\
''PHP is a widely-used open source general-purpose scripting language that is especially suited for web development and can be embedded into HTML'' \cite{php}. We use PHP to inform the google cloud messaging service if there's a new image on our web server.


\textbf{FTP - File Transfer Protocol} \\
‘’The File Transfer Protocol is a standard network protocol used to transfer computer files from one host to another host over a TCP-based network, such as the Internet''. In our system the arduino uses the ftp protocol to send the image over a passive ftp connection to the web server.

\textbf{Android} \\
Android is a operating system for mobile devices. It is based on the linux kernel and so android is open source. For an easy development we used different external libraries such as google play services or picasso. Picasso allows simple and easy image loading in our application. Google play services is needed for the google cloud messaging service and we also had to include the gcm library.


\textbf{Google Cloud Messaging Service} \\
One key aspect of the android app is the google cloud messaging service. We used this technology because it is easy to implement and we needed a notification system. Otherwise our app had to poll all the time and this would be very energy inefficient. With the gcm we have now a push notification whenever our web server receives a new image.

\begin{figure}[H]
	\centering
	\includegraphics[width=0.9\textwidth]{images/usecase.png}
	\caption[Google Cloud Messaging]{This image shows the structure of the google cloud messaging.}
	\label{gcm}
\end{figure}

As you can see in figure \ref{gcm} the android device sends his senderID and his applicationID to the google cloud messaging service server for registration. In the second step, the gcm server sends a registrationID back to the device. With this registrationID, the mobile phone can be registered from our server. Therefor the device sends the registrationID to the server (point three in the picture \ref{gcm}). The registrationID will be stored in a database for quicker access. This id is important, because whenever an event occours and a push notification is needed, our server sends a message in a string to the google cloud messaging server along with the registrationID (a). Every android device which is registered in the gcm server under the same registrationID will be notified (b). So it is possible to send the same notification and message to various android devices. This will make the spreading of the taken image easier. For example if more than one person lives in the house where motion was detected or the picture can be send to more than one facility manager.
\section[User interface - Benjamin Kugler]{User interface}
\section[Conclusion - Team]{Conclusion}


\end{absaetze}

% Literaturverzeichnis
% http://ctan.mackichan.com/macros/latex/exptl/biblatex/doc/biblatex.pdf - 2.1.1 Regular Types
%\addcontentsline{toc}{section}{Literatur}
%\bibliographystyle{gerabbrv}
%\bibliographystyle{alpha}
\bibliographystyle{plain}
\bibliography{bib/references}

\renewcommand{\thepage}{}
\end{document}


